\documentclass[11pt]{article}
\usepackage[T1]{fontenc}
\usepackage{textcomp}
\usepackage[margin=0.75in]{geometry}
\usepackage{amsmath,amssymb,amsthm}
\usepackage{array,booktabs}
\usepackage{hyperref}
\setlength{\parindent}{0pt}
\newtheorem{theorem}{Theorem}
\theoremstyle{definition}
\newtheorem{definition}{Definition}
\title{Flow Proof Artifact: Nat-mul}
\author{Generated by \texttt{flow doc proof}}
\date{Flow Proof Book}
\begin{document}
\maketitle
Peano definitions for multiplication on natural numbers.\\[0.5em]
\textbf{Source.} peano — https://en.wikipedia.org/wiki/Peano\_axioms\\[0.5em]
\bigskip
\noindent{\large \textbf{Definition 1.} \textit{Zero times anything is zero} \hfill the natural numbers \cdot multiplication \cdot zero is the left annihilator}\par
\smallskip\noindent\textit{Coordinate: the natural numbers \cdot multiplication \cdot zero is the left annihilator}\par
\smallskip\noindent\textit{Source: peano}\par
\medskip
\noindent\textbf{Goal.} Zero times anything is zero.  (0 * 7 = 0)\par
\noindent$\displaystyle \forall m \in \mathbb{N}\quad 0 \cdot m = 0$\par
\medskip
\noindent
\begin{tabular}{@{} >{\raggedright\arraybackslash}p{0.06\textwidth} >{\raggedright\arraybackslash}p{0.47\textwidth} >{\raggedleft\arraybackslash}p{0.41\textwidth} @{}}
\textbf{\#} & \textbf{Proof} & \textbf{Mathematics} \\
\midrule
\textbf{1.} & We stipulate zero is the left annihilator for multiplication on the natural numbers — this is a definition, not a derived fact. & \\[0.45em]
\textbf{2.} & This follows directly from the definition: 0 times m equals 0. Hence proven. & $\displaystyle 0 \cdot m = 0$ \\[0.45em]
\bottomrule
\end{tabular}
\label{thm:Nat-multiplication-zero_is_the_left_annihilator}
\par\medskip\hrule
\bigskip
\noindent{\large \textbf{Definition 2.} \textit{Multiplying by a successor adds another copy} \hfill the natural numbers \cdot multiplication \cdot successor on the right distributes}\par
\smallskip\noindent\textit{Coordinate: the natural numbers \cdot multiplication \cdot successor on the right distributes}\par
\smallskip\noindent\textit{Source: peano}\par
\medskip
\noindent\textbf{Goal.} Multiplying by a successor adds another copy.  (3 * succ(2) = 3*2 + 3)\par
\noindent$\displaystyle \forall n \in \mathbb{N} \forall m \in \mathbb{N}\quad n * \mathrm{succ}(m) = n \cdot m + n$\par
\medskip
\noindent
\begin{tabular}{@{} >{\raggedright\arraybackslash}p{0.06\textwidth} >{\raggedright\arraybackslash}p{0.47\textwidth} >{\raggedleft\arraybackslash}p{0.41\textwidth} @{}}
\textbf{\#} & \textbf{Proof} & \textbf{Mathematics} \\
\midrule
\textbf{1.} & We stipulate successor on the right distributes for multiplication on the natural numbers — this is a definition, not a derived fact. & \\[0.45em]
\textbf{2.} & This follows directly from the definition: n times the successor of m equals n times m plus n. Hence proven. & $\displaystyle n * \mathrm{succ}(m) = n \cdot m + n$ \\[0.45em]
\bottomrule
\end{tabular}
\label{thm:Nat-multiplication-successor_on_the_right_distributes}
\par\medskip\hrule
\bigskip
\noindent{\large \textbf{Definition 3.} \textit{Multiplying anything by zero on the right gives zero} \hfill the natural numbers \cdot multiplication \cdot zero is the right annihilator}\par
\smallskip\noindent\textit{Coordinate: the natural numbers \cdot multiplication \cdot zero is the right annihilator}\par
\smallskip\noindent\textit{Source: peano}\par
\medskip
\noindent\textbf{Goal.} Multiplying anything by zero on the right gives zero.  (7 * 0 = 0)\par
\noindent$\displaystyle \forall n \in \mathbb{N}\quad n \cdot 0 = 0$\par
\medskip
\noindent
\begin{tabular}{@{} >{\raggedright\arraybackslash}p{0.06\textwidth} >{\raggedright\arraybackslash}p{0.47\textwidth} >{\raggedleft\arraybackslash}p{0.41\textwidth} @{}}
\textbf{\#} & \textbf{Proof} & \textbf{Mathematics} \\
\midrule
\textbf{1.} & We stipulate zero is the right annihilator for multiplication on the natural numbers — this is a definition, not a derived fact. & \\[0.45em]
\textbf{2.} & This follows directly from the definition: n times 0 equals 0. Hence proven. & $\displaystyle n \cdot 0 = 0$ \\[0.45em]
\bottomrule
\end{tabular}
\label{thm:Nat-multiplication-zero_is_the_right_annihilator}
\par\medskip\hrule
\end{document}
